\chapter{Содержание курсового проекта}

Работа включает пять тем --- построение кривых элементов теоретического чертежа, построение масштаба Бонжана, расчет остойчивости на больших углах крена, расчет непотопляемости и расчет продольного спуска.

\section{Кривые элементов теоретического чертежа}

Кривые элементов теоретического чертежа — зависимости элементов плавучести и начальной остойчивости судна от осадки.
Элементы плавучести и начальной остойчивости с их обозначениями приведены в таблице \ref{tab:elements}.

Элементы погруженного объема --- водоизмещение, абсцисса и аппликата центра величины — используются при определении осадки судна на ровный киль.

Элементы площади ватерлинии --- ее площадь, абсцисса центра тяжести, моменты инерции площади относительно продольной и поперечный осей --- используются для вычисления метацентрических характеристик корпуса судна --- поперечного и продольного метацентрических радиусов и аппликаты метацентра.

Дополнительно могут быть также построены зависимости коэффициентов полноты от осадки.

\begin{table}[h]
    \caption{Элементы теоретического чертежа}
    \label{tab:elements}
    \centering
    \small
    \begin{tabular}{ c  l  }
        Обозначение & Наименование \\ [1.5ex]
        \hline \\
        $V$ & Объемное водоизмещение \\ [1.5ex]
        $x_c$ & Абсцисса центра величины (центра подводного объема) \\ [1.5ex]
        $z_c$ & Аппликата центра величины \\ [1.5ex]
        $S$ & Площадь ватерлинии \\ [1.5ex]
        $x_f$ & Абсцисса центра площади ватерлинии \\ [1.5ex]
        $I_x$ & Момент инерции площади ватерлинии относительно оси $X$ \\ [1.5ex]
        $I_y$ & Момент инерции площади ватерлинии относительно оси $Y$ \\ [1.5ex]
        $I_{yf}$ & Момент инерции площади ватерлинии \\ [1.5ex]
         & относительно поперечной центральной оси \\ [1.5ex]
        $r$ & Поперечный метацентрический радиус\\ [1.5ex]
        $R$ & Продольный метацентрический радиус\\ [1.5ex]
        $z_m$ & Возвышение поперечного метацентра\\ [1.5ex]
    \end{tabular}
\end{table}

В первом задании курсового проекта (или домашней работы), используя исходные данные, необходимо выполнить расчет и построить кривые элементов теоретического чертежа.
Исходные данные включают размерную таблицу ординат шпангоутов, главные размерения, шпацию, расстояние между ватерлиниями.

Рассмотрим необходимые формулы и последовательность выполнения расчетов.

Площадь ватерлинии вычисляется для каждой из ватерлиний кроме нулевой (ВЛ0), так как в этом случае судно в принципе не погружено в воду и понятие ватерлинии не имеет смысла.
Формула для вычисления площади ватерлинии в общем виде

\begin{equation}
    S = \int_{x_{\text{К}}}^{x_{\text{Н}}} (y_{\text{пр}} - y_{\text{лб}}) dx.
\end{equation}

В случае симметричных относительно ДП бортов (посадка судна без крена) $y_{\text{пр}} = - y_{\text{лб}}$ выражение для площади ватерлинии принимает вид

\begin{equation}
    S = 2 \int_{x_{\text{К}}}^{x_{\text{Н}}} y dx.
\end{equation}

Применяя правило трапеций с учетом коэффициентов, получим

\begin{equation}
    S = 2 \Delta L (K_{\text{Ф}} y_{\text{Н}} + y_{\text{Н}+1} + \dots + y_{\text{К}-1} + K_{\text{А}} y_{\text{К}}).
\end{equation}

Абсцисса центра тяжести площади ватерлинии определяется по общеизвестной формуле для плоских фигур как отношение статического момента площади ватерлинии к ее площади

\begin{equation}
    x_f = M_y / S.
\end{equation}

Статический момент площади ватерлинии относительно поперечной оси $Y$ определяется по формуле

\begin{equation}
    M_y = \int_{x_{\text{К}}}^{x_{\text{Н}}} y \cdot x dx.
\end{equation}

Применяя правило трапеций с учетом коэффициентов, получим

\begin{equation}
    \begin{split}
    M_y = 2 \Delta L^2 [K_{\text{Ф}} y_{\text{Н}} (5 - i_{\text{Н}}) + y_{\text{Н}+1}(5 - i_{\text{Н}+1})+ \dots \\
    + y_{\text{К}-1}(5-i_{\text{К}-1}) + K_{\text{А}} y_{\text{К}} (5 - i_{\text{К}})].
    \end{split}
\end{equation}

Величинами $i_j$ с соответствующими индексами обозначается номер шпангоута.
При этом произведение $\Delta L (5 - i_j)$ --- расстояние от шпангоута до начала координат, совпадающего с мидель–шпангоутом.
В случае разбивки не на 10, а на 20 шпангоутов, вместо 5 в формуле нужно использовать 10 (то есть номер мидель-шпангоута).

Схема вычисления площади ватерлинии и статического момента площади ватерлинии приведена на рисунке \ref{tab:waterline-elements}.
В приведенной схеме ватерлиния притыкается к ДП в носу между шпангоутами 0 и 1, а в корме между шпангоутами 9 и 10.

% Схема расчета площади ватерлинии и статического момента площади ватерлинии от осадки в табличном редакторе
\begin{figure}
    \footnotesize
    \begin{tikzpicture}
        \matrix (mat) [table,
        row 1/.style={nodes={fill=gray!10,align=center,font=\segoe}},
        column 1/.style={nodes={fill=gray!10,align=center,font=\segoe}},
        column 1/.append style={nodes={text width=0.5cm}},
        column 2/.append style={nodes={text width=1.0cm}},
        column 3/.append style={nodes={text width=1.0cm}},
        column 4/.append style={nodes={text width=4.0cm}},
        column 5/.append style={nodes={text width=4.0cm}}
        ] {
        & A  & B & C & D \\
        1 & 0 & ---   & ---      & --- \\
        2 & 1 & $y_1$ & =$K_{\text{Ф}}$*B2 & =$K_{\text{Ф}}$*B2*(5-A2)\\
        3 & 2 & $y_2$ & =B3 & =B3*(5-A3)\\
        4 & 3 & $y_3$ & =B4 & =B4*(5-A4)\\
        5 & 4 & $y_4$ & =B5 & =B5*(5-A5)\\
        6 & 5 & $y_5$ & =B6 & =B6*(5-A6)\\
        7 & 6 & $y_6$ & =B7 & =B7*(5-A7)\\
        8 & 7 & $y_7$ & =B8 & =B8*(5-A8)\\
        9 & 8 & $y_8$ & =B9 & =B9*(5-A9)\\
        10 & 9 & $y_9$ & =$K_{\text{А}}$*B10 & =$K_{\text{А}}$*B10*(5-A10)\\
        11 & 10 & --- & --- & --- \\
        12 & --- & --- & =2*$\Delta L$*СУММ(C2:C10) & =2*$\Delta L$\string^2*СУММ(D2:D10) \\
        };
    \end{tikzpicture}
    \caption{Схема расчета площади ватерлинии и статического момента площади ватерлинии от осадки в табличном редакторе}
    \caption*{
        \textit{
        В столбце A приведены номера шпангоутов.\\
        В столбце B приведены ординаты шпангоутов на рассматриваемой ватерлинии.
        Если рассматриваемая ватерлиния и шпангоут не пересекается, в соответствующей ячейке ставится прочерк.\\
        В столбце C приведены значения ординат ватерлинии с учетом корректирующих коэффициентов учета истинного притыкания ватерлинии к ДП.
        В последней строке столбца C вычисляется значение площади ватерлинии.
        $ \Delta L$ в формуле --- расстояние между шпангоутами, вычисленное ранее.\\
        В столбце D приведены значения ординат ватерлинии, умноженные на безразмерные абсциссы шпангоутов.
        В последней строке столбца D вычисляется значение статического момента площади ватерлинии относительно оси $Y$.}
        }
    \label{tab:waterline-elements}
\end{figure}

Продольная центральная ось в случае симметрии бортов ватерлинии совпадает с координатной осью $X$.
Момент инерции площади ватерлинии относительно координатной оси $X$ вычисляется по формуле

\begin{equation}
    I_x = \frac{2}{3} \int_{x_{\text{К}}}^{x_{\text{Н}}} y^3 dx.
\end{equation}

Применяя правило трапеции с учетом коэффициентов, получим

\begin{equation}
    I_x = \frac{2}{3} \Delta L (K_{\text{Ф}} y_{\text{Н}}^3 + y_{\text{Н}+1}^3 + \dots + y_{\text{К}-1}^3 + K_{\text{А}} y_{\text{К}}^3).
\end{equation}

Поперечная центральная ось ватерлинии в общем случае не совпадает с поперечной координатной осью $Y$, поэтому вычисление момента инерции относительно поперечной центральной оси выполняется в два этапа.
Сначала определяется момент инерции площади ватерлинии относительно координатной оси $Y$

\begin{equation}
    I_y = 2 \int_{x_{\text{К}}}^{x_{\text{Н}}} y x^2 dx.
\end{equation}

Применяя правило трапеций с учетом коэффициентов, получим

\begin{equation}
    \begin{split}
    I_y = 2 \Delta L^3 [K_{\text{Ф}} y_{\text{Н}} (5 - i_{\text{Н}})^2 + y_{\text{Н}+1} (5 - i_{\text{Н}+1})^2 + \dots \\
    + y_{\text{К}-1} (5 - i_{\text{К}-1})^2 + K_{\text{А}} y_{\text{К}} (5 - i_{\text{К}})^2].
    \end{split}
\end{equation}

А затем вычисляется собственно момент инерции площади ватерлинии относительно поперечной центральной оси

\begin{equation}
    I_{yf} = I_y - S \cdot x_f^2.
\end{equation}

Схема вычисления моментов инерции площади ватерлинии относительно продольной и поперечной осей на рисунке \ref{tab:inertia-elements}.

% Схема расчета моментов инерции площади ватерлинии относительно продольной и поперечной осей в табличном редакторе
\begin{figure}
    \centering
    \footnotesize
    \begin{tikzpicture}
        \matrix (mat) [table,
        row 1/.style={nodes={fill=gray!10,align=center,font=\segoe}},
        column 1/.style={nodes={fill=gray!10,align=center,font=\segoe}},
        column 1/.append style={nodes={text width=0.5cm}},
        column 2/.append style={nodes={text width=1.2cm}},
        column 3/.append style={nodes={text width=1.2cm}},
        column 4/.append style={nodes={text width=4cm}},
        column 5/.append style={nodes={text width=4cm}}
        ] {
        & A  & B & C & D \\
        1 & 0 & ---   & ---      & --- \\
        2 & 1 & $y_1$ & =$K_{\text{Ф}}$*B2\string^3 & =$K_{\text{Ф}}$*B2*(5-A2)\string^2\\
        3 & 2 & $y_2$ & =B3\string^3 & =B3*(5-A3)\string^2\\
        4 & 3 & $y_3$ & =B4\string^3 & =B4*(5-A4)\string^2\\
        5 & 4 & $y_4$ & =B5\string^3 & =B5*(5-A5)\string^2\\
        6 & 5 & $y_5$ & =B6\string^3 & =B6*(5-A6)\string^2\\
        7 & 6 & $y_6$ & =B7\string^3 & =B7*(5-A7)\string^2\\
        8 & 7 & $y_7$ & =B8\string^3 & =B8*(5-A8)\string^2\\
        9 & 8 & $y_8$ & =B9\string^3 & =B9*(5-A9)\string^2\\
        10 & 9 & $y_9$ & =$K_{\text{А}}$*B10\string^3 & =$K_{\text{А}}$*B10*(5-A10)\string^2\\
        11 & 10 & --- & --- & --- \\
        12 & --- & --- & =2/3*$\Delta L$*СУММ(C2:C10) & = 2*$\Delta L$\string^3*СУММ(D2:D10) \\
        };
    \end{tikzpicture}
    \caption{Схема расчета моментов инерции площади ватерлинии относительно продольной и поперечной осей в табличном редакторе}
    \caption*{\itshape{
        В столбце A приведены номера шпангоутов.\\
        В столбце B приведены ординаты шпангоутов на рассматриваемой ватерлинии.
        Если рассматриваемая ватерлиния и шпангоут не пересекается, в соответствующей ячейке ставится прочерк.\\
        В столбце C приведены значения ординат ватерлинии с учетом корректирующих коэффициентов учета истинного притыкания ватерлинии к ДП в третьей степени.
        В последней строке столбца C вычисляется значение момента инерции площади ватерлинии относительно продольной оси.\\
        В столбце D приведены значения ординат ватерлинии, умноженные на квадрат безразмерных абсцисс шпангоутов.
        В последней строке столбца D вычисляется значение момента инерции площади ватерлинии относительно оси $Y$.}}
    \label{tab:inertia-elements}
\end{figure}  

Расчеты перечисленных элементов выполняются для каждой ватерлинии.
В результате получаем зависимости $S$, $x_f$, $M_y$, $I_x$, $I_y$ и $I_{yf}$ от осадки \ref{tab:waterline-elements-results}.

% Элементы площади ватерлинии в зависимости от осадки в табличном редакторе
\begin{figure}
    \centering
    \footnotesize
    \begin{tikzpicture}
        \matrix (mat) [table,
        row 1/.style={nodes={fill=gray!10,align=center,font=\segoe}},
        column 1/.style={nodes={fill=gray!10,align=center,font=\segoe}},
        column 1/.append style={nodes={text width=0.5cm}},
        column 2/.append style={nodes={text width=1.0cm}},
        column 3/.append style={nodes={text width=1.5cm}},
        column 4/.append style={nodes={text width=1.0cm}},
        column 5/.append style={nodes={text width=1.0cm}},
        column 6/.append style={nodes={text width=1.0cm}},
        column 7/.append style={nodes={text width=1.0cm}},
        column 8/.append style={nodes={text width=1.0cm}},
        column 9/.append style={nodes={text width=1.0cm}}
        ] {
        & A  & B & C & D & E & F & G & H \\
        1 & 0 & =$\Delta T$*A1 & $S_0$ & $x_{f0}$ & $M_{y0}$ & $I_{x0}$ & $I_{y0}$ & $I_{yf0}$ \\
        2 & 1 & =$\Delta T$*A2 & $S_1$ & $x_{f1}$ & $M_{y1}$ & $I_{x1}$ & $I_{y1}$ & $I_{yf1}$ \\
        3 & 2 & =$\Delta T$*A3 & $S_2$ & $x_{f2}$ & $M_{y2}$ & $I_{x2}$ & $I_{y2}$ & $I_{yf2}$ \\
        4 & 3 & =$\Delta T$*A4 & $S_3$ & $x_{f3}$ & $M_{y3}$ & $I_{x3}$ & $I_{y3}$ & $I_{yf3}$ \\
        5 & 4 & =$\Delta T$*A5 & $S_4$ & $x_{f4}$ & $M_{y4}$ & $I_{x4}$ & $I_{y4}$ & $I_{yf4}$ \\
        6 & 5 & =$\Delta T$*A6 & $S_5$ & $x_{f5}$ & $M_{y5}$ & $I_{x5}$ & $I_{y5}$ & $I_{yf5}$ \\
        7 & 6 & =$\Delta T$*A7 & $S_6$ & $x_{f6}$ & $M_{y6}$ & $I_{x6}$ & $I_{y6}$ & $I_{yf6}$ \\
        8 & 7 & =$\Delta T$*A8 & $S_7$ & $x_{f7}$ & $M_{y7}$ & $I_{x7}$ & $I_{y7}$ & $I_{yf7}$ \\
        };
    \end{tikzpicture}
    \caption{Элементы площади ватерлинии в зависимости от осадки в табличном редакторе}
    \caption*{\textit{
        В столбце B вычисляется осадка на соответствующей ватерлинии.
        $\Delta T$ в формулах --- расстояние между ватерлиниями, вычисленное ранее.
    }}
    \label{tab:waterline-elements-results}
\end{figure}

Объемное водоизмещение вычисляется по формуле

\begin{equation}
    V(z) = \int_0^z S dz.
\end{equation}

Для вычисления абсциссы и аппликаты центра величины необходимо вычислить две вспомогательные величины.
Статический момент объема относительно плоскости мидель–шпангоута и статический момент объема относительно основной плоскости.

Статический момент объема относительно плоскости мидель–шпангоута вычисляется по формуле

\begin{equation}
    M_{yz} = \int_0^z S x_f dz = \int_0^z M_y dz.
\end{equation}

Статический момент объема относительно основной плоскости вычисляется по формуле

\begin{equation}
    M_{xy} = \int_0^z S z dz.
\end{equation}

Абсцисса и аппликата центра величины вычисляются с использованием полученных величин

\begin{equation}
    \begin{split}
    x_c = \frac{M_{yz}}{V}, \\
    z_c = \frac{M_{xy}}{V}.
    \end{split}
\end{equation}

Схема расчета объемного водоизмещения, абсциссы и аппликаты центра величины приведена на рисунке \ref{tab:volume-elements}.

% Схема расчета объемного водоизмещения и статических моментов объема в табличном редакторе
\begin{figure}
    \centering
    \footnotesize
    \begin{tikzpicture}
        \matrix (mat) [table,
        row 1/.style={nodes={fill=gray!10,align=center,font=\segoe}},
        column 1/.style={nodes={fill=gray!10,align=center,font=\segoe}},
        column 1/.append style={nodes={text width=0.5cm}},
        column 2/.append style={nodes={text width=0.7cm}},
        column 3/.append style={nodes={text width=0.7cm}},
        column 4/.append style={nodes={text width=0.7cm}},
        column 5/.append style={nodes={text width=4.0cm}},
        column 6/.append style={nodes={text width=4.0cm}}
        ] {
        & A  & B & C & D & E  \\
        1 & 0 & $S_0$ & $M_{y0}$ & 0 & 0 \\
        2 & 1 & $S_1$ & $M_{y1}$ & =D1+0,5*(B1+B2)*$\Delta T$ & =E1+0,5*(C1+C2)*$\Delta T$ \\
        3 & 2 & $S_2$ & $M_{y2}$ & =D2+0,5*(B2+B3)*$\Delta T$ & =E2+0,5*(C2+C3)*$\Delta T$ \\
        4 & 3 & $S_3$ & $M_{y3}$ & =D3+0,5*(B3+B4)*$\Delta T$ & =E3+0,5*(C3+C4)*$\Delta T$ \\
        5 & 4 & $S_4$ & $M_{y4}$ & =D4+0,5*(B4+B5)*$\Delta T$ & =E4+0,5*(C4+C5)*$\Delta T$ \\
        6 & 5 & $S_5$ & $M_{y5}$ & =D5+0,5*(B5+B6)*$\Delta T$ & =E5+0,5*(C5+C6)*$\Delta T$ \\
        7 & 6 & $S_6$ & $M_{y6}$ & =D6+0,5*(B6+B7)*$\Delta T$ & =E6+0,5*(C6+C7)*$\Delta T$ \\
        8 & 7 & $S_7$ & $M_{y7}$ & =D7+0,5*(B7+B8)*$\Delta T$ & =E7+0,5*(C7+C8)*$\Delta T$ \\
        };
    \end{tikzpicture}

    \begin{tikzpicture}
        \matrix (mat) [table,
        row 1/.style={nodes={fill=gray!10,align=center,font=\segoe}},
        column 1/.style={nodes={fill=gray!10,align=center,font=\segoe}},
        column 1/.append style={nodes={text width=0.5cm}},
        column 2/.append style={nodes={text width=0.7cm}},
        column 3/.append style={nodes={text width=0.7cm}},
        column 4/.append style={nodes={text width=0.7cm}},
        column 5/.append style={nodes={text width=5.5cm}},
        column 6/.append style={nodes={text width=1.0cm}},
        column 7/.append style={nodes={text width=1.0cm}}
        ] {
        & A  & B & C & F & G & H  \\
        1 & 0 & $S_0$ & $M_{y0}$ & 0 & --- & --- \\
        2 & 1 & $S_1$ & $M_{y1}$ & =F1+0,5*(A1*B1+A2*B2)*$\Delta T$*$\Delta T$ & =E2/D2 & =F2/D2 \\
        3 & 2 & $S_2$ & $M_{y2}$ & =F2+0,5*(A2*B2+A3*B3)*$\Delta T$*$\Delta T$ & =E3/D3 & =F3/D3 \\
        4 & 3 & $S_3$ & $M_{y3}$ & =F3+0,5*(A3*B3+A4*B4)*$\Delta T$*$\Delta T$ & =E4/D4 & =F4/D4 \\
        5 & 4 & $S_4$ & $M_{y4}$ & =F4+0,5*(A4*B4+A5*B5)*$\Delta T$*$\Delta T$ & =E5/D5 & =F5/D5 \\
        6 & 5 & $S_5$ & $M_{y5}$ & =F5+0,5*(A5*B5+A6*B6)*$\Delta T$*$\Delta T$ & =E6/D6 & =F6/D6 \\
        7 & 6 & $S_6$ & $M_{y6}$ & =F6+0,5*(A6*B6+A7*B7)*$\Delta T$*$\Delta T$ & =E7/D7 & =F7/D7 \\
        8 & 7 & $S_7$ & $M_{y7}$ & =F7+0,5*(A7*B7+A8*B8)*$\Delta T$*$\Delta T$ & =E8/D8 & =F8/D8 \\
        };
    \end{tikzpicture}

    \caption{Схема расчета объемного водоизмещения и статических моментов объема в табличном редакторе}
    \caption*{\itshape{
        В столбце D вычисляется объемное водоизмещение для каждой ватерлинии (интегрирование с переменным верхним пределом).\\
        В столбце E вычисляется статический момент объема относительно плоскости мидель–шпангоута для каждой ватерлинии.\\
        В столбце F вычисляется статический момент объема относительно основной плоскости для каждой ватерлинии.\\
        В столбцах G и H вычисляются абсцисса и аппликата центра величины соответственно для каждой ватерлинии.
    }}
    \label{tab:volume-elements}
\end{figure}

Метацентрические характеристик корпуса судна вычисляются следующим образом.

Поперечный метацентрический радиус

\begin{equation}
    r = \frac{I_x}{V}.
\end{equation}

Продольный метацентрический радиус

\begin{equation}
    R = \frac{I_{yf}}{V}.
\end{equation}

Возвышение (аппликата) поперечного метацентра

\begin{equation}
    z_m = z_C + r.
\end{equation}

Указанные величины вычисляются для каждой ватерлинии, кроме нулевой, так как в этом случае объемное водоизмещение равно нулю, а ватерлиния отсутствует, и понятие метацентрических радиусов не имеют смысла.
Схема расчета метацентрических радиусов и возвышения метацентра приведена на рисунке \ref{tab:metacentric-radii}.

% Схема расчета метацентрических радиусов и возвышения метацентра в табличном редакторе
\begin{figure}
    \centering
    \footnotesize
    \begin{tikzpicture}
        \matrix (mat) [table,
        row 1/.style={nodes={fill=gray!10,align=center,font=\segoe}},
        column 1/.style={nodes={fill=gray!10,align=center,font=\segoe}},
        column 1/.append style={nodes={text width=0.5cm}},
        column 2/.append style={nodes={text width=0.7cm}},
        column 3/.append style={nodes={text width=1.0cm}},
        column 4/.append style={nodes={text width=1.0cm}},
        column 5/.append style={nodes={text width=1.0cm}},
        column 6/.append style={nodes={text width=1.0cm}},
        column 7/.append style={nodes={text width=1.5cm}},
        column 8/.append style={nodes={text width=1.5cm}},
        column 9/.append style={nodes={text width=1.5cm}}
        ] {
        & A  & B & C & D & E & F & G & H  \\
        1 & 0 & 0 & $I_{x0}$ & $I_{yf0}$ & $z_{c0}$ & --- & --- & --- \\
        2 & 1 & $V_1$ & $I_{x1}$ & $I_{yf1}$ & $z_{c1}$ & =C2/B2 & =D2/B2 & =E2+F2 \\
        3 & 2 & $V_2$ & $I_{x2}$ & $I_{yf2}$ & $z_{c2}$ & =C3/B3 & =D3/B3 & =E3+F3 \\
        4 & 3 & $V_3$ & $I_{x3}$ & $I_{yf3}$ & $z_{c3}$ & =C4/B4 & =D4/B4 & =E4+F4 \\
        5 & 4 & $V_4$ & $I_{x4}$ & $I_{yf4}$ & $z_{c4}$ & =C5/B5 & =D5/B5 & =E5+F5 \\
        6 & 5 & $V_5$ & $I_{x5}$ & $I_{yf5}$ & $z_{c5}$ & =C6/B6 & =D6/B6 & =E6+F6 \\
        7 & 6 & $V_6$ & $I_{x6}$ & $I_{yf6}$ & $z_{c6}$ & =C7/B7 & =D7/B7 & =E7+F7 \\
        8 & 7 & $V_7$ & $I_{x7}$ & $I_{yf7}$ & $z_{c7}$ & =C8/B8 & =D8/B8 & =E8+F8 \\
        };
    \end{tikzpicture}
    \caption{Схема расчета метацентрических радиусов и возвышения метацентра в табличном редакторе}
    \caption*{\itshape{
        В столбце F вычисляется поперечный метацентрический радиус для каждой ватерлинии.\\
        В столбце G вычисляется продольный метацентрический радиус для каждой ватерлинии.\\
        В столбце H вычисляется возвышение (аппликата) метацентра для каждой ватерлинии.
    }}
    \label{tab:metacentric-radii}
\end{figure}

Как сказано выше, кривые элементов теоретического чертежа используются для вычисления характеристик остойчивости при различной осадке судна, а значит, оформление кривых должно быть выполнено в удобной для этого форме.

В отчете по курсовой работе диаграмма с кривыми элементов теоретического чертежа строится на листе формата А3 в альбомной ориентации.
Масштаб диаграмм, построенных в табличном редакторе (\textit{Microsoft Excel} или \textit{LibreOffice Calc}), крайне сложно настроить поэтому кривые элементов теоретического чертежа следует строить в чертежной программе (например, \textit{Kompas-3D} или \textit{NanoCAD}).
Если кривые вычерчиваются вручную, то используйте миллиметровую бумагу.

Масштаб по вертикальной оси (по осадке) должен быть выбран стандартным из перечня ЕСКД (1:50, 1:100, 1:200 и т.п.).
На вертикальной оси кроме делений в метрах должны быть нанесены метки, соответствующие расчетным ватерлиниям с указанием их обозначения (ВЛ1, ВЛ2 и т.п.).

Горизонтальная ось --- ось, по которой откладываются значения элементов теоретического чертежа — следует градуировать в сантиметрах и их долях.

Масштаб кривых следует выбирать стандартным из перечня ЕСКД для каждого элемента теоретического чертежа.
Выбранный масштаб следует указывать непосредственно рядом с кривой или в технических требованиях рядом с диаграммой.
Выбирать масштабы следует так, чтобы максимально полно использовать поле диаграммы, вместе с тем так, чтобы кривые не перекрывали друг друга.

Масштаб кривых абсцисс центра тяжести площади ватерлинии и центра величины должен быть одинаковым, также как и масштаб поперечного метацентрического радиуса, аппликаты центра величины и возвышения метацентра.

Пример оформления кривых элементов теоретического чертежа приведен в Приложении.

Проконтролировать правильность выполненных расчетов при построении кривых элементов теоретического чертежа можно по следующим признакам:

--- все кривые имеют плавный характер, без изломов и крутых перегибов;

--- в точке пересечения кривых $x_f$ и $x_c$ последняя должна иметь экстремум;

--- величина $z_c$ при любой осадке должна быть в пределах от ее половины до двух третей;

--- отношение $I_{yf}/I_x$ и $R/r$ должны иметь порядок $(L/B)^2$;

--- продольный метацентрический радиус $R$ для КВЛ должен иметь порядок длины $L$.

\section{Масштаб Бонжана}

\textbf{Масштаб Бонжана} --- диаграмма, состоящая из кривых зависимости площади шпангоута от осадки, вычерченных на шпангоутах.
Основное назначение масштаба Бонжана --- определение водоизмещения и абсциссы центра величины при наличии у судна дифферента.

Кроме кривых площади шпангоутов на масштабе Бонжана наносят линию основной плоскости, вертикальные линии теоретических шпангоутов, контуры штевней, бортовую линию верхней палубы и линию палубы в ДП.
Все линии, кроме кривых площади шпангоутов строятся в двух масштабах --- один для вертикальной оси осадок и второй для горизонтальной оси по длине судна.
Масштаб кривых площади шпангоутов выбирается стандартным из перечня ЕСКД так, чтобы диаграммой было удобно пользоваться.

В этом задании курсового проекта (или домашней работы) необходимо построить масштаб Бонжана, используя исходные данные, и определить водоизмещение и абсциссу центра величины судна при наличии у него дифферента.

Итак, рассмотрим порядок расчета и построения масштаба Бонжана.
Площадь шпангоута определяется по формуле

\begin{equation}
    \Omega(z) = \int_0^z y dz.
\end{equation}

Из рисунка \ref{fig:station-area-1} видно, что всю площадь шпангоута можно условно разделить на три части: от нулевой до верхней ватерлинии $\Omega_{\text{р}}$, от верхней ватерлинии до уровня палубной линии у борта $\Delta \Omega$ и от уровня палубной линии у борта до уровня палубы в ДП $\Delta \Omega_f$.
Для удобства рассмотрим вычисление площадей этих участков раздельно.

% К расчету площади шпангоута
\begin{figure}[h!]
    \centerline{\incfig{figure_09 - 1}}
    \caption{К расчету площади шпангоута}
    \caption*{\textit{
        Слева схематично приведено шпангоутное сечение с тремя частями площади.\\
        Справа --- кривая погруженной части этого шпангоута в зависимости от осадки.
    }}
    \label{fig:station-area-1}
\end{figure}

Площадь шпангоута от нулевой до верхней ватерлинии $\Omega_{\text{р}}$ вычисляется с использованием правила трапеции по формуле

\begin{equation}
    \Omega(z) = \int_0^z 2 y dz \approx \Delta T \sum_0^z (y_j + y_{j+1}).
\end{equation}

Площадь от верхней ватерлинии до уровня палубной линии у борта $\Delta \Omega$ вычисляется по формуле с учетом предположения прямобортности как площадь трапеции

\begin{equation}
    \Delta \Omega = (z_{\text{П}} - z_m) (y_{\text{П}} - y_{m}),
\end{equation}
где $y_m$, $z_m$ --- ордината и аппликата верхней расчетной ватерлинии, $y_{\text{П}}$, $z_{\text{П}}$ --- ордината и аппликата палубной линии у борта.

Площадь от уровня палубной линии у борта до палубы в ДП $\Delta \Omega_f$ представляет собой площадь сегмента круга и вычисляется по формуле

\begin{equation}
    \Delta \Omega_f = 4/3 \cdot y_{\text{П}} \cdot f = \frac{4}{75} y_{\text{П}}^2.
\end{equation}

% Схема расчета зависимости площади шпангоута от осадки в табличном редакторе
\begin{figure}
    \centering
    \footnotesize
    \begin{tikzpicture}
        \matrix (mat) [table,
        row 1/.style={nodes={fill=gray!10,align=center,font=\segoe}},
        column 1/.style={nodes={fill=gray!10,align=center,font=\segoe}},
        column 1/.append style={nodes={text width=0.5cm}},
        column 2/.append style={nodes={text width=1.5cm}},
        column 3/.append style={nodes={text width=1.5cm}},
        column 4/.append style={nodes={text width=1.5cm}},
        column 5/.append style={nodes={text width=4.5cm}}
        ] {
        & A  & B & C & D \\
        1 & 0 & 0 & $y_0^*$ & 0 \\
        2 & 1 & =$\Delta T$*A2 & $y_1$ & =D1+$\Delta T$*(C1+C2) \\
        3 & 2 & =$\Delta T$*A3 & $y_2$ & =D2+$\Delta T$*(C2+C3) \\
        4 & 3 & =$\Delta T$*A4 & $y_3$ & =D3+$\Delta T$*(C3+C4) \\
        5 & 4 & =$\Delta T$*A5 & $y_4$ & =D4+$\Delta T$*(C4+C5) \\
        6 & 5 & =$\Delta T$*A6 & $y_5$ & =D5+$\Delta T$*(C5+C6) \\
        7 & 6 & =$\Delta T$*A7 & $y_6$ & =D6+$\Delta T$*(C6+C7) \\
        8 & 7 & =$\Delta T$*A8 & $y_7$ & =D7+$\Delta T$*(C7+C8) \\
        9 & --- & $z_{\text{П}}$ & $y_{\text{П}}$ & =D8+(B9-B8)*(C8+C9) \\
        10& --- & --- & --- & =D9+4/75*C9\string^2 \\
        };
    \end{tikzpicture}
    \caption{Схема расчета зависимости площади шпангоута от осадки в табличном редакторе}
    \caption*{\textit{
        В столбце B вычисляется осадка на соответствующем уровне.
        В девятой строке приведено возвышение палубной линии у борта.\\
        В столбце C приведены ординаты шпангоута на соответствующем уровне.
        В десятой строке приведено значение ординаты палубной линии у борта.\\
        В столбце D вычисляется площадь шпангоута.
        В девятой строке к площади шпангоута по верхнюю ватерлинию добавляется площадь от верхней ватерлинии до палубной линии у борта.
        В десятой строке к площади шпангоута добавляется поправка для учета погиби палубы.
        }}
    \label{tab:station-area-bonjean}
\end{figure}

Зависимость площади шпангоута от осадки вычисляется для каждого шпангоута.

При построении масштаба Бонжана следует выполнять некоторые условия:

-- масштаб Бонжана вычерчивается с использованием CAD-программы, например \textit{Kompas-3D} или \textit{NanoCAD}, или вручную на миллиметровой бумаге;

-- горизонтальная ось изображает основную плоскость, по которой в натурном для судна масштабе отмечаются положения шпангоутов вертикальными линиями;

-- вертикальная ось изображает осадки, по которой в натурном для судна масштабе отмечаются ватерлинии горизонтальными линиями;

-- на носовом и кормовом перпендикулярах (нулевой и десятый шпангоуты) наносятся шкалы осадки в метрах;

-- кривые площадей шпангоутов наносятся на каждом шпангоуте в масштабе, выбранном из стандартных масштабов ЕСКД, так, чтобы диаграммой было удобно пользоваться.

% Пример масштаба Бонжана
\begin{figure}[h!]
    \centerline{\incfig{figure_10 - 1}}
    \caption{Пример масштаба Бонжана}
    \caption*{\textit{
        На каждом шпангоуте построена кривая зависимости площади погруженной части этого шпангоута от осадки. \\
        Проведены горизонтальные линии, соответствующие ватерлиниям.\\
        Наклонная прямая, построенная между осадкой на носовом и кормовом перпендикулярах, соответствует ватерлинии судна при дифференте. \\
        Величина площади погруженной части на каждом шпангоуте определяется для точки пересечения этой ватерлинии с вертикальной линией шпангоута.
    }}
    \label{fig:bonjean-diagram}
\end{figure}


Ниже представлена схема определения водоизмещения и абсциссы центра тяжести судна при посадке с дифферентом с использованием масштаба Бонжана.

В связи с тем, что масштаб Бонжана построен в неравных горизонтальной и вертикальной осях (масштаб по осадке больше масштаба по длине) использовать способ задания посадки судна по осадке на миделе и углу дифферента невозможно.
В первую очередь следует определить осадки на носовом и кормовом перпендикулярах

\begin{equation}
    \begin{split}
    T_f = T + \psi (\frac{L}{2} - x_f), \\
    T_a = T - \psi (\frac{L}{2} + x_f),
    \end{split}
\end{equation}
где $T$ --- осадка судна на ровный киль, м;
$x_f$ --- абсцисса центра тяжести площади ватерлинии при осадке $T$.
В формулах выше используется значение угла дифферента в радианах вместо тангенса этого угла (при условии достаточной его малости).

На масштабе Бонжана на носовом и кормовом перпендикулярах откладываются соответствующие значения осадок, через которые проводится наклонная прямая ватерлиния.

В точках пересечения построенной ватерлинии с вертикальными линиями шпангоутов для каждого из них следует снять значения площади с кривых $\Omega(z)$.
С использованием этих данных выполняется вычисление водоизмещения и абсциссы центра тяжести

\begin{equation}
    \begin{split}
    V = \int_{L_{\text{К}}}^{L_{\text{Н}}} \Omega dx, \\
    x_C = \frac{1}{V} \int_{L_{\text{К}}}^{L_{\text{Н}}} \Omega x dx.
    \end{split}
\end{equation}

Применяя правило трапеции, преобразуем интегральные выражения в суммы

\begin{equation}
    \begin{split}
    V = \Delta L (K_{\text{А}} \cdot \Omega_{\text{К}} + \dots + \Omega_{\text{Н}+1} + K_{\text{Ф}} \cdot \Omega_{\text{Н}}), \\
    x_C = \Delta L / V (K_{\text{А}} \cdot \Omega_{\text{К}} \cdot x_{\text{К}} + \dots + \Omega_{\text{Н}+1} \cdot x_{\text{Н}+1} + K_{\text{Ф}} \cdot \Omega_{\text{Н}} \cdot x_{\text{Н}}).
    \end{split}
\end{equation}

Коэффициенты $K_{\text{А}}$ и $K_{\text{Ф}}$ принимаются равными этим коэффициентам для ватерлиний, ближайших к наклонной ватерлинии в корме и носу соответственно.

Схема расчета водоизмещения и абсциссы центра величины судна, сидящего с дифферентом, приведена на рисунке \ref{fig:volume-elements-bonjean}.

% Схема расчета водоизмещения и абсциссы центра величины судна по масштабу Бонжана в табличном редакторе
\begin{figure}
    \centering
    \footnotesize
    \begin{tikzpicture}
        \matrix (mat) [table,
        row 1/.style={nodes={fill=gray!10,align=center,font=\segoe}},
        column 1/.style={nodes={fill=gray!10,align=center,font=\segoe}},
        column 1/.append style={nodes={text width=0.5cm}},
        column 2/.append style={nodes={text width=0.7cm}},
        column 3/.append style={nodes={text width=0.7cm}},
        column 4/.append style={nodes={text width=3.5cm}},
        column 5/.append style={nodes={text width=4.5cm}}
        ] {
        & A  & B & C & D \\
        1 & 0 & --- & --- & --- \\
        2 & 1 & $\Omega_1$ & =$K_{\text{Ф}}$*B2 & =$K_{\text{Ф}}$*B2*(5 - A2) \\
        3 & 2 & $\Omega_2$ & =B3 & =B3*(5 - A3) \\
        4 & 3 & $\Omega_3$ & =B4 & =B4*(5 - A4) \\
        5 & 4 & $\Omega_4$ & =B5 & =B5*(5 - A5) \\
        6 & 5 & $\Omega_5$ & =B6 & =B6*(5 - A6) \\
        7 & 6 & $\Omega_6$ & =B7 & =B7*(5 - A7) \\
        8 & 7 & $\Omega_7$ & =B8 & =B8*(5 - A8) \\
        9 & 8 & $\Omega_8$ & =B9 & =B9*(5 - A9) \\
        10 & 9 & $\Omega_9$ & =$K_{\text{А}}$*B10 & =$K_{\text{А}}$*B10*(5 - A10) \\
        11 & 10 & --- & --- & --- \\
        12 & --- & --- & =$\Delta L$*СУММ(C2:C10) & =$\Delta L$\string^2*СУММ(D2:D10)/C12 \\
        };
    \end{tikzpicture}
    \caption{Схема расчета водоизмещения и абсциссы центра величины судна по масштабу Бонжана в табличном редакторе}
    \caption*{\textit{Значения в столбце B снимаются с кривых площади шпангоутов на масштабе Бонжана. Коэффициенты $K_{\text{А}}$ и $K_{\text{Ф}}$ принимаются равными этим коэффициентам для ватерлиний, ближайших к наклонной ватерлинии в корме и носу соответственно.
    В схеме крайними шпангоутами приведены первый и девятый.}}
    \label{fig:volume-elements-bonjean}
\end{figure}

\section{Остойчивость на больших углах крена}

Для оценки остойчивости на больших углах крена используются диаграммы статической и динамической остойчивости, а также интерполяционные кривые.
Настоящий раздел курсовой работы посвящен их расчету и построению.

Плечо статической остойчивости при угле крена $\Theta$ вычисляется по формуле

\begin{equation}
    l_{\Theta} = l_{\text{ф}} - a \cdot \sin \Theta,
\end{equation}
где
$l_{\text{ф}}$ --- плечо остойчивости формы, м;
$a = z_g - z_c$ --- расстояние от центра тяжести судна до центра величины, м.

Плечо динамической остойчивости при угле крена $\Theta$ вычисляется по формуле

\begin{equation}
    d = \int_0^{\Theta} l_{\Theta} d\Theta .
\end{equation}

Аппликата центра тяжести $z_g$ задается в исходных данных к курсовому проекту, аппликата центра величины вычисляется с использованием элементов теоретического чертежа (Задание 1 курсовой работы).
Поэтому для расчета диаграмм статической и динамической остойчивости в первую очередь требуется рассчитать величины плеч остойчивости формы $l_{\text{ф}}$ для различных углов крена.

Для расчета плеч остойчивости формы используется вспомогательный теоретический чертеж.
Для заданного водоизмещения проводят последовательное накренение на выбранные заранее углы крена в диапазоне от 0\textdegree до 90\textdegree.
Шаг по углу крена в курсовой работу может приниматься равным 15\textdegree.

Из основного курса статики известно, что статическое накренение судна равнообъемно, то есть водоизмещение судна при изменении угла крена не меняется.
В связи с тем, что форма корпуса судна сложна, решение задачи равнообъемного накренения не тривиально.

В компьютерных программах по теории корабля решение задачи равнообъемных наклонений может выполняться последовательными приближениями (решением оптимизационной задачи), или интерполированием между заранее просчитанными элементами теоретического чертежа при различных углах крена.
Подобные методы, хотя и могут быть применены без специальных программных средств, сложны к реализации с использованием табличных редакторов.
Поэтому в курсовой работе используется графо–аналитический способ Крылова–Дарньи.

Метод основан на проведении вспомогательной ватерлинии наклонения, отсекающей приближенно постоянный объем, с уточнением положения равнообъемной ватерлинии определением расстояния между ней и вспомогательной ватерлинией.

Вспомогательная ватерлиния каждый раз проводится через центр тяжести предыдущей (по углу крена) равнообъемной ватерлинии. Способ легко осуществляется при выполнении расчетов статики вручную.

% Построение равнообъемной ватерлинии методом Крылова–Дарньи
\begin{figure}[h!]
    \centerline{\incfig{figure_11 - 1}}
    \caption{Построение равнообъемной ватерлинии методом Крылова–Дарньи}
    \label{fig:kryloff-darny}
\end{figure}

Плечо остойчивости формы при накренении на угол $\Theta$ определяется по формуле

\begin{equation}
    l_{\text{ф}} = y_{\Theta} \cos \Theta + (z_{\Theta} - z_c) \sin \Theta,
    \label{eq:form_stability_arm}
\end{equation}
где $y_{\Theta}$, $z_{\Theta}$ --- ордината и аппликата центра величины равнообъемной ватерлинии при угле крена $\Theta$;
$z_c$ --- аппликата центра величины судна в положении без крена.

Координаты центра величины при накренении на угол $\Theta$ рассчитываются по формулам

\begin{equation}
    \begin{split}
    y_{\Theta} = \int_0^{\Theta} r_{\Theta} \cos \Theta d\Theta, \\
    z_{\Theta} -z_c = \int_0^{\Theta} r_{\Theta} \sin \Theta d\Theta .
    \end{split}
\end{equation}

В формулах используется поперечный метацентрический радиус судна, накрененного на угол $\Theta$, определяемый по формуле

\begin{equation}
    r_{\Theta} = \frac{I_{x\Theta}}{V},
\end{equation}
где $I_{x\Theta}$ --- момент инерции площади накрененной равнообъемной ватерлинии при угле крена $\Theta$ относительно ее центральной продольной оси.

При описании расчета кривых элементов теоретического чертежа было отмечено, что центральная продольная ось совпадает (обычно) с продольной координатной осью $X$, лежащей в ДП.
Однако, в случае накренения судна говорить о симметричности его ватерлиний уже нельзя, поэтому центральная ось не совпадает с координатной продольной осью.
В связи с этим, для вычисления момента инерции площади накрененной равнообъемной ватерлинии относительно ее центральной продольной оси, необходимо сначала определить координаты центра величины этой ватерлинии, а затем вычислить момент инерции относительно координатной оси $X$ и применить теорему параллельных осей.

Последовательность выполнения расчетов:

-- расчет значений момента инерции и метацентрического радиуса для равнообъемных наклонных ватерлиний;

-- расчет и построение кривых остойчивости формы;

-- построение интерполяционных кривых;

-- расчет и построение диаграммы статической и динамической остойчивости.

Рассмотрим подробно каждый из этих этапов.

\subsection{Расчет момента инерции площади ватерлинии и метацентрического радиуса при накренении}

Способ Крылова–Дарньи требует чередования вычислений и графических построений на теоретическом чертеже корпуса.

Порядок выполнения расчетов:

Для начальной ватерлинии (при нулевом угле крена) вычисляется момент инерции ватерлинии относительно центральной продольной оси и метацентрический радиус.

Через центр тяжести начальной ватерлинии (начальная точка $F_0$ находится на пересечении ДП и следа ватерлинии) сплошной линией проводится вспомогательная ватерлиния.
Шаг по углу крена в курсовой работе допускается принимать равным 15\textdegree.
С теоретического чертежа снимаются значения ординат шпангоутов в точках их пересечения со вспомогательной ватерлинией.
Так как правый и левый борта при накренении несимметричны, снимать значения ординат нужно на оба борта.

Правило знаков: ординаты, направленные в сторону входящего в воду борта, положительны.
Ординаты, направленные в противоположную от входящего в воду борта сторону, отрицательные.

По снятым с чертежа ординатам вычисляется площадь и статический момент площади ватерлинии

\begin{equation}
    \begin{split}
    S_{\Theta} = \int_{L_{\text{К}}}^{L_{\text{Н}}} (y_{\text{пр}} - y_{\text{лб}}) dx, \\
    M_{x\Theta} = \int_{L_{\text{К}}}^{L_{\text{Н}}} (y_{\text{пр}}^2 - y_{\text{лб}}^2) dx.
    \end{split}
\end{equation}

Ордината центра тяжести площади ватерлинии накрененного судна определяется по формуле

\begin{equation}
    y_{f \Theta} = \frac{M_{x\Theta}}{S_{\Theta}}.
\end{equation}

Далее вычисляется момент инерции площади ватерлинии

\begin{equation}
    I_{x \Theta} = I_{x \Theta}^\prime - \Delta I_{x \Theta}.
\end{equation}

Здесь

\begin{equation}
    I_{x \Theta}^\prime = 1/3 \int_{L_{\text{К}}}^{L_{\text{Н}}} (y_{\text{пр}}^3 - y_{\text{лб}}^3) dx,
\end{equation}

\begin{equation}
    \Delta I_{x \Theta} = S_{\Theta} \cdot y_{f \Theta}^2.
\end{equation}

Поперечный метацентрический радиус вычисляется по формуле

\begin{equation}
    r_{\Theta} = \frac{I_{x \Theta}}{V}.
\end{equation}

Учет истинной длины ватерлинии выполняется с использованием коэффициентов $K_{\text{А}}$ и $K_{\text{Ф}}$, которые определяются для посадки судна на ровный киль при заданном для этого раздела курсовой работы водоизмещении.
Значения коэффициентов используются неизменными для всех углов накренения.

% Построение равнообъемной ватерлинии методом Крылова–Дарньи
\begin{figure}
    \centering
    \footnotesize
    \begin{tikzpicture}
        \matrix (mat) [table,
        row 1/.style={nodes={fill=gray!10,align=center,font=\segoe}},
        column 1/.style={nodes={fill=gray!10,align=center,font=\segoe}},
        column 1/.append style={nodes={text width=0.5cm}},
        column 2/.append style={nodes={text width=0.7cm}},
        column 3/.append style={nodes={text width=0.8cm}},
        column 4/.append style={nodes={text width=0.8cm}},
        column 5/.append style={nodes={text width=3.6cm}},
        column 6/.append style={nodes={text width=3.6cm}}
        ] {
        & A  & B & C & D & E \\
        1 & 0 & $y_{{\text{пр}}0}$ & $y_{\text{лб}0}$ & =$K_{\text{Ф}}$*(B1-C1) & =$K_{\text{Ф}}$*(B1\string^2-C1\string^2) \\
        2 & 1 & $y_{{\text{пр}}1}$ & $y_{\text{лб}1}$ & =(B2-C2) & =(B2\string^2-C2\string^2) \\
        3 & 2 & $y_{{\text{пр}}2}$ & $y_{\text{лб}2}$ & =(B3-C3) & =(B3\string^2-C3\string^2) \\
        4 & 3 & $y_{{\text{пр}}3}$ & $y_{\text{лб}3}$ & =(B4-C4) & =(B4\string^2-C4\string^2) \\
        5 & 4 & $y_{{\text{пр}}4}$ & $y_{\text{лб}4}$ & =(B5-C5) & =(B5\string^2-C5\string^2) \\
        6 & 5 & $y_{{\text{пр}}5}$ & $y_{\text{лб}5}$ & =(B6-C6) & =(B6\string^2-C6\string^2) \\
        7 & 6 & $y_{{\text{пр}}6}$ & $y_{\text{лб}6}$ & =(B7-C7) & =(B7\string^2-C7\string^2) \\
        8 & 7 & $y_{{\text{пр}}7}$ & $y_{\text{лб}7}$ & =(B8-C8) & =(B8\string^2-C8\string^2) \\
        9 & 8 & $y_{{\text{пр}}8}$ & $y_{\text{лб}8}$ & =(B9-C9) & =(B9\string^2-C9\string^2) \\
        10 & 9 & $y_{{\text{пр}}9}$ & $y_{\text{лб}9}$ & =(B10-C10) & =(B10\string^2-C10\string^2) \\
        11 & 10 & $y_{{\text{пр}}10}$ & $y_{\text{лб}10}$ & =$K_{\text{А}}$*(B11-C11) & =$K_{\text{А}}$*(B11\string^2-C11\string^2) \\
        12 & --- & --- & --- & =$\Delta L$*СУММ(D1:D11) & =$\Delta L$*СУММ(E1:E11)/2 \\
        };
    \end{tikzpicture}
    \begin{tikzpicture}
        \matrix (mat) [table,
        row 1/.style={nodes={fill=gray!10,align=center,font=\segoe}},
        column 1/.style={nodes={fill=gray!10,align=center,font=\segoe}},
        column 1/.append style={nodes={text width=0.5cm}},
        column 2/.append style={nodes={text width=0.7cm}},
        column 3/.append style={nodes={text width=0.8cm}},
        column 4/.append style={nodes={text width=0.8cm}},
        column 5/.append style={nodes={text width=6.0cm}}
        ] {
        & A  & B & C  & F \\
        1 & 0 & $y_{{\text{пр}}0}$ & $y_{\text{лб}0}$ & =$K_{\text{Ф}}$*(B1\string^3 - C1\string^3) \\
        2 & 1 & $y_{{\text{пр}}1}$ & $y_{\text{лб}1}$ & =(B2\string^3 - C2\string^3) \\
        3 & 2 & $y_{{\text{пр}}2}$ & $y_{\text{лб}2}$ & =(B3\string^3 - C3\string^3) \\
        4 & 3 & $y_{{\text{пр}}3}$ & $y_{\text{лб}3}$ & =(B4\string^3 - C4\string^3) \\
        5 & 4 & $y_{{\text{пр}}4}$ & $y_{\text{лб}4}$ & =(B5\string^3 - C5\string^3) \\
        6 & 5 & $y_{{\text{пр}}5}$ & $y_{\text{лб}5}$ & =(B6\string^3 - C6\string^3) \\
        7 & 6 & $y_{{\text{пр}}6}$ & $y_{\text{лб}6}$ & =(B7\string^3 - C7\string^3) \\
        8 & 7 & $y_{{\text{пр}}7}$ & $y_{\text{лб}7}$ & =(B8\string^3 - C8\string^3) \\
        9 & 8 & $y_{{\text{пр}}8}$ & $y_{\text{лб}8}$ & =(B9\string^3 - C9\string^3) \\
        10 & 9 & $y_{{\text{пр}}9}$ & $y_{\text{лб}9}$& =(B10\string^3 - C10\string^3) \\
        11 & 10 & $y_{{\text{пр}}10}$ & $y_{\text{лб}10}$ & =$K_{\text{А}}$*(B11\string^3 - C11\string^3) \\
        12 & --- & --- & --- & =$\Delta L$*СУММ(F1:F11)/3 \\
        };
    \end{tikzpicture}
    \caption{Схема расчета равнообъемных наклонений в табличном редакторе}
    \caption*{\textit{Здесь $y_{\text{пр}}$ и $y_{\text{лб}}$ --- ординаты правого и левого борта соответственно, снятые с чертежа по вспомогательной ватерлинии.
    В курсовой работе принято накренение на правый борт.\\
    В столбце D вычисляется площадь ватерлинии, в столбце E --- статический момент площади, в столбце F --- момент инерции площади.}}
    \label{tab:kryloff-method-II}
\end{figure}

Рассчитанные по вспомогательной ватерлинии величины позволяют построить равнообъемную наклонную ватерлинию.
Для этого вдоль по ватерлинии с предыдущим углом наклона от точки $F_{\Theta-1}$ откладывают отрезок $1/2y_{f\Theta}$.
Через полученную точку проводят равнообъемную наклонную ватерлинию.
Положительное значение $1/2y_{f\Theta}$ откладывается в сторону входящего в воду борта, отрицательное --- в противоположную.
Пример приведен на рисунке \ref{fig:kryloff-darny-large}.

Для определения центра тяжести площади равнообъемной наклонной ватерлинии из центра тяжести предыдущей равнообъемной ватерлинии вдоль по вспомогательной ватерлинии откладывается значение $y_{f\Theta}$.
Пересечение нормали к вспомогательной ватерлинии, проведенной через полученную точку, с равнообъемной ватерлинией есть точка $F_{\Theta}$ --- центр тяжести площади наклонной равнообъемной ватерлинии.

% Построение равнообъемной ватерлинии методом Крылова–Дарньи
\begin{figure}[h!]
    \centerline{\incfig{figure_11 - 2}}
    \caption{Построение равнообъемной ватерлинии методом Крылова–Дарньи}
    \label{fig:kryloff-darny-large}
\end{figure}

Далее повторяются уже описанные выше шаги при накренении на следующий угол крена.
Таким образом, в курсовой работе проводятся построения и вычисления сначала для угла крена 15\textdegree.
Затем, использую полученные данные, для угла крена 30\textdegree, затем 45\textdegree и 60\textdegree.
При желании можно выполнить дополнительные расчеты для углов крена 75\textdegree и 90\textdegree.

После вычисления метацентрического радиуса для всех углов крена строится график $r=f(\Theta)$.
Пример графика приведен на рисунке \ref{fig:metacentric-radius-heel}.

Следует обратить внимание на некоторые особенности таких функций.
В связи с симметрией корпуса относительно ДП, касательная к функции при отсутствии крена (0\textdegree) горизонтальна.
А в случае резкого изменения площади ватерлинии (выход скулы их воды или вход в нее палубы) на кривой может появиться местный максимум.

\begin{figure}[h!]
    \centerline{\incfig{figure_12 - 1}}
    \caption{Зависимость метацентрического радиуса от угла крена}
    \caption*{\textit{В курсовой работе выполняется расчет зависимости метацентрического радиуса для двух значений водоизмещения, соответствующих осадкам по расчетную ватерлинию ГВЛ-2 и ГВЛ+1}}
    \label{fig:metacentric-radius-heel}
\end{figure}

\subsection{Расчет и построение кривых остойчивости формы}

Морские суда эксплуатируются при различных случаях нагрузки, характеризующихся величиной дедвейта и положением его центра тяжести.
Для каждого из этих случаев диаграмма статической остойчивости своя.
Чтобы облегчить построение диаграммы статической остойчивости для любого случая нагрузки используют интерполяционные кривые плеч остойчивости формы.

Плечи остойчивости формы зависят только от формы погруженного объема корпуса судна, но не от распределения нагрузки, поэтому могут быть посчитаны заранее при проектировании судна.

Для построения интерполяционных кривых задаются несколькими равномерно распределенными значениями водоизмещения, охватывающими полный диапазон возможных значений водоизмещения судна --- от водоизмещения порожнем до водоизмещения полностью погруженного судна (верхняя палуба под водой).
Для этих значений водоизмещения определяют зависимость плеча остойчивости формы от угла крена по формуле \ref{eq:form_stability_arm}.
Для водоизмещения погруженного судна при всех углах крена плечи остойчивости формы равны нулю, так как действующая ватерлиния отсутствует.

В курсовом проекте следует выполнить расчет плеч остойчивости формы для трех значений водоизмещения:
водоизмещения порожнем и двух промежуточных значений водоизмещения до полностью погруженного судна.
За осадку при водоизмещении порожнем можно принять осадку по расчетной ватерлинии на две ниже КВЛ.
Полученные кривые $l_{\text{ф}}=f(\Theta)$ построить на одной диаграмме. 
Для контроля правильности построения кривых можно использовать свойство формулы, по которой выполнен расчет плеч остойчивости формы

\begin{equation}
    \frac{d l_{\text{ф}}}{d\Theta} l_{\Theta = 0} = r_0.
\end{equation}

Графически это значит, что если на графике при угле крена $\Theta = 1\,\text{рад}$ отложить значение метацентрического радиуса при осадке без крена и провести линию через эту точку и начало координат, то кривая плеча остойчивости формы будет касательной к такому отрезку.

Для того чтобы по полученной диаграмме зависимости плеч остойчивости формы от угла крена при нескольких значениях водоизмещения построить интерполяционные кривые, необходимо посчитать значение водоизмещения, при котором судно находится полностью под водой.
Для этого нужно проинтегрировать площади шпангоутов по верхнюю палубу, полученные с использованием масштаба Бонжана

\begin{equation}
    V_{\text{П}} = \int_{L_{\text{К}}}^{L_{\text{Н}}} \Omega_{\text{П}} dx.
\end{equation}

При расчете объема с применением правила трапеции для крайних шпангоутов для учета объемов оконечностей нужно использовать поправочные коэффициенты $K_{\text{А}}$ и $K_{\text{Ф}}$ (в курсовой работе можно принять равным значениям этих коэффициентов для самой верхней ватерлинии).

% Зависимость плеча остойчивости формы от угла крена для различных водоизмещений
\begin{figure}[h!]
    \centerline{\incfig{figure_13 - 1}}
    \caption{Зависимость плеча остойчивости формы от угла крена для различных водоизмещений}
    \caption*{\textit{Кривые построены для двух значений водоизмещения, соответствующих осадкам по расчетную ватерлинию ГВЛ-2 и ГВЛ+1}}
    \label{fig:form-stability-arm}
\end{figure}

\subsection{Построение интерполяционных кривых}

Интерполяционные кривые плеч остойчивости формы --- представление кривых остойчивости формы в других осях.
Если кривые плеч остойчивости формы строятся для нескольких значений водоизмещения в осях угол крена --- плечо, то интерполяционные кривые строятся в осях водоизмещение --- плечо.
При этом конкретная кривая соответствует определенному значению угла крена.

В курсовой работе плечи остойчивости формы определяются для четырех значений угла крена (15\textdegree, 30\textdegree, 45\textdegree и 60\textdegree), поэтому кривых будет четыре.
По горизонтальной оси откладывают значения водоизмещения, для которых выполнен расчет плеч остойчивости формы, а на вертикальные линии, проведенные через них, откладывают значения плеч для четырех значений угла крена.
При значении водоизмещения полного погружения судна все кривые должны сойтись в одной точке, лежащей на горизонтальной оси.
Пример интерполяционных кривых остойчивости формы приведен на рисунке \ref{fig:interpolation-curves-stability}.

По интерполяционным кривым достаточно легко определить значения плеч остойчивости формы для любого значения водоизмещения и угла крена.

% Интерполяционные кривые остойчивости формы
\begin{figure}[h!]
    \centerline{\incfig{figure_14 - 1}}
    \caption{Интерполяционные кривые остойчивости формы}
    \label{fig:interpolation-curves-stability}
\end{figure}

\subsection{Расчет и построение диаграмм статической и динамической остойчивости}

Диаграмма статической остойчивости --- графическое представление зависимости плеча остойчивости $l_{\Theta}$ от угла крена.
В предыдущем разделе приведен процесс вычисления плеча остойчивости формы, величина $a = z_g-z_c$ включает значение аппликаты центра тяжести (в курсовом проекте можно принять равным $0,7 H$, где $H$ --- высота борта).

Схема расчета диаграмм статической и динамической остойчивости приведена на рисунке \ref{fig:stability-curve-calculation}.

% Схема расчета диаграмм статической и динамической остойчивости в табличном редакторе
\begin{figure}
    \centering
    \footnotesize
    \begin{tikzpicture}
        \matrix (mat) [table,
        row 1/.style={nodes={fill=gray!10,align=center,font=\segoe}},
        column 1/.style={nodes={fill=gray!10,align=center,font=\segoe}},
        column 1/.append style={nodes={text width=0.5cm}},
        column 2/.append style={nodes={text width=1.5cm}},
        column 3/.append style={nodes={text width=0.7cm}},
        column 4/.append style={nodes={text width=1.5cm}},
        column 5/.append style={nodes={text width=3.0cm}},
        column 6/.append style={nodes={text width=3.0cm}}
        ] {
        & A  & B & C & D & E \\
        1 & =SIN(0) & $l_{\text{ф}0}$ & =B1-$a$*A1 & 0 & =$\Delta \Theta$/2 * D1 \\
        2 & =SIN(15) & $l_{\text{ф}15}$ & =B2-$a$*A2 & =D1+(C1+C2) & =$\Delta \Theta$/2 * D2 \\
        3 & =SIN(30) & $l_{\text{ф}30}$ & =B3-$a$*A3 & =D2+(C2+C3) & =$\Delta \Theta$/2 * D3 \\
        4 & =SIN(45) & $l_{\text{ф}45}$ & =B4-$a$*A4 & =D3+(C3+C4) & =$\Delta \Theta$/2 * D4 \\
        5 & =SIN(60) & $l_{\text{ф}60}$ & =B5-$a$*A5 & =D4+(C4+C5) & =$\Delta \Theta$/2 * D5 \\
        };
    \end{tikzpicture}
    \caption{Схема расчета диаграмм статической и динамической остойчивости в табличном редакторе}
    \caption*{\textit{Значения в столбце $C$ соответствуют значениям плеч статической остойчивости.
    Значения в столбце $E$ --- значениям плеч динамической остойчивости.
    В столбце $A$ значение угла в градусах приведено условно --- функция $SIN$ принимает значение угла только в радианах.}}
    \label{fig:stability-curve-calculation}
\end{figure}

Правильность расчета и построения диаграммы статической остойчивости можно проверить, используя определение начальной метацентрической высоты, которая является касательной к кривой плеч статической остойчивости.
Если отложить значение начальной метацентрической высоты на перпендикуляре, отложенном от горизонтальной оси при угле крена, равном 1 радиану, и провести прямую через отложенную точку и начало координат, то эта прямая должна быть касательной к диаграмме статической остойчивости при отсутствии крена.

Плечо динамической остойчивости определяется как интеграл с переменным верхним пределом зависимости плеча статической остойчивости от угла крена, поэтому между этими диаграммами существуют интегральные соотношения:

-- при углах крена, соответствующих точкам пересечения диаграммы статической остойчивости с горизонтальной осью (плечо равно нулю), диаграмма динамической остойчивости имеет экстремум (максимум или минимум);

-- при угле крена, соответствующем максимуму диаграммы статической остойчивости, диаграмма динамической остойчивости имеет перегиб.

\begin{figure}[h!]
    \centerline{\incfig{figure_14 - 2}}
    \caption{Диаграммы статической и динамической остойчивости}
    \label{fig:stability-curves}
\end{figure}

\section{Непотопляемость судна}

Раздел включает построение кривой предельной длины отсеков, а также расчет посадки и начальной остойчивости поврежденного судна.

\subsection{Построение кривой предельной длины отсеков}

Предельная длина отсека --- расстояние между поперечными водонепроницаемыми переборками, при котором еще выполняются заданные условия непотопляемости (аварийной плавучести и аварийной остойчивости).
Предельная длина отсека зависит от положения этого отсека по длине судна.
Зависимость предельной длины отсека от положения по длине судна называется кривой предельной длины отсека.

Требование аварийной плавучести заключается в том, что при затоплении одного или нескольких отсеков судно должно сидеть не глубже предельной линии погружения, назначаемой в соответствии с правилами классификационного общества.

Кривую предельной длины отсеков строят с использованием масштаба Бонжана для заданных значений водоизмещения, а также положения центра тяжести судна по длине и высоте ($x_g$ и $z_g$).
По кривым элементов теоретического чертежа определяются: средняя осадка, координаты центра величины, метацентрические радиусы, площадь ватерлинии, абсцисса центра тяжести площади ватерлинии, центральные моменты инерции площади ватерлинии и начальный дифферент.

На масштаб Бонжана наносится предельная линия погружения (в соответствии с конкретными требованиями классификационного общества).
К этой кривой проводится несколько касательных (желательно, нечетное количество), представляющих собой следы предельных ватерлиний на диаметральной плоскости.

Для проведения крайних предельных ватерлиний на носовом и кормовом перпендикулярах откладывается расстояние $z_0 = 1{,}2(2 T_{\text{КВЛ}} - H)$.
Из полученных точек строятся касательные к линии предельного погружения (в примере линия предельного погружения совпадает с линией палубы у борта).
Если провести касательную не получается, линия проводится то точки пересечения штевня с предельной линией погружения.

Средняя предельная ватерлиния должна проводиться параллельно основной плоскости.
Углы между остальными касательными ватерлиниями должны быть примерно одинаковы.

% К расчету кривой предельных объемов
\begin{figure}[h!]
    \centerline{\incfig{figure_15 - 1}}
    \caption{К расчету кривой предельных объемов}
    \caption*{\textit{Серыми линиями на масштабе Бонжана показано пять предельных ватерлиний, необходимых для построения кривой предельных объемов.
    В курсовой работе рекомендуется выполнять расчет для семи предельных ватерлиний.}}
    \label{fig:bonjean-volumes}
\end{figure}

Для каждой из предельных ватерлиний с масштаба Бонжана снимаются значения площади погруженной части шпангоутов.
По этим значениям интегрированием определяется водоизмещение и статические моменты погруженного объема относительно мидель–шпангоута $V_{i-i}$, где $i$ --- номер предельной ватерлинии.

Объем отсека, при затоплении которого судно погрузится по предельную ватерлинию определяется как разница рассчитанного и начального значения водоизмещения

\begin{equation}
    v_i = V_{i-i} - V_0.
\end{equation}

Абсцисса центра тяжести объема $v_i$ определяется как отношение разницы статических моментов к объему

\begin{equation}
    x_i = \frac{M_{i-i} - M_0}{v_i},
\end{equation}
где $M_0 = V_0 \cdot x_c$ --- статический момент водоизмещения судна до повреждения.

Схема расчета объема отсеков, при затоплении которых судно погрузится до предельной ватерлинии приведено на рисунке \ref{fig:full-volume}.

% Схема расчета водоизмещения и статического момента при посадке по предельную ватерлинию в табличном редакторе
\begin{figure}
    \centering
    \footnotesize
    \begin{tikzpicture}
        \matrix (mat) [table,
        row 1/.style={nodes={fill=gray!10,align=center,font=\segoe}},
        column 1/.style={nodes={fill=gray!10,align=center,font=\segoe}},
        column 1/.append style={nodes={text width=0.5cm}},
        column 2/.append style={nodes={text width=0.7cm}},
        column 3/.append style={nodes={text width=1.3cm}},
        column 4/.append style={nodes={text width=3.8cm}},
        column 5/.append style={nodes={text width=4.2cm}}
        ] {
        & A  & B & C & D \\
        1 & 0 & --- & --- & --- \\
        2 & 1 & $\Omega_1$ & =$K_{\text{Ф}}$*B2 & =$K_{\text{Ф}}$*B2*(5-A2) \\
        3 & 2 & $\Omega_2$ & =B3 & =B3*(5-A3) \\
        4 & 3 & $\Omega_3$ & =B4 & =B4*(5-A4) \\
        5 & 4 & $\Omega_4$ & =B5 & =B5*(5-A5) \\
        6 & 5 & $\Omega_5$ & =B6 & =B6*(5-A6) \\
        7 & 6 & $\Omega_6$ & =B7 & =B7*(5-A7) \\
        8 & 7 & $\Omega_7$ & =B8 & =B8*(5-A8) \\
        9 & 8 & $\Omega_8$ & =B9 & =B9*(5-A9) \\
        10 & 9 & $\Omega_9$ & =B10 & =B10*(5-A10) \\
        11 & 10 & $\Omega_10$ & =$K_{\text{А}}$*B11 & =$K_{\text{А}}$*B11*(5-A11) \\
        12 & --- & --- & =$\Delta L$*СУММ(C2:C11) & =$\Delta L$\string^2*СУММ(D2:D11) \\
        };
    \end{tikzpicture}
    \caption{Схема расчета водоизмещения и статического момента при посадке по предельную ватерлинию в табличном редакторе}
    \caption*{\textit{В примере в качестве крайних носового и кормового шпангоутов приведены первый и десятый.}}
    \label{fig:full-volume}
\end{figure}

Коэффициенты, приведенные в схеме расчета, используются для учета объема оконечностей и принимаются равными значениям этих коэффициентов, рассчитанных для ближайших ватерлиний.

Далее необходимо построить кривую предельных объемов отсеков.
Для этого на горизонтальной прямой отмечается мидель-шпангоут, от которого в масштабе откладываются значения абсцисс центра тяжести объема влившейся воды.
От полученных точек вверх в масштабе откладывается значение соответствующего объема.
Точки соединяют плавной кривой (кривой предельных объемов отсеков).

На этой же диаграмме вычерчивают строевую по шпангоутам по предельную линию погружения.
Необходимые для этого значения площадей шпангоутов снимают с масштаба Бонжана.

% К расчету кривой предельных длин отсеков
\begin{figure}[h!]
    \centerline{\incfig{figure_15 - 2}}
    \caption{К расчету кривой предельных длин отсеков}
    \caption*{\textit{На верхней диаграмме нанесены объемы отсеков в соответствующих им абсциссах центров тяжести.
    В нижней части нанесена строевая по шпангоутам по линию предельного погружения (в своем масштабе) и интегральная кривая (в масштабе объемов).\\
    $1$ --- кривая предельных объемов отсеков.\\
    $2$ --- строевая по шпангоутам по предельную линию погружения.\\
    $3$ --- интегральная кривая $V_X$.
    }}
    \label{fig:max-volume-curve}
\end{figure}
% Пример определения предельных длин отсеков

Для вычисления предельных длин отсеков с использованием кривой предельных объемов отсеков необходимо рассчитать и нанести на диаграмму интегральную кривую от строевой по шпангоутам --- кривую, представляющую собой в каждой точке по длине значение объемного водоизмещения части судна в нос от этой точки.
Вычисляется интегральная кривая как интеграл с переменным верхним пределом от строевой по шпангоутам

\begin{equation}
    V_X = \int_{L_{\text{К}}}^{L_{\text{Н}}} \Omega_{\text{П}} dx,
\end{equation}
где $V_X$ --- объем по предельную ватерлинию от носа до поперечного сечения с абсциссой $x$.

Схема расчета интегральной кривой по строевой по шпангоутам по правилу трапеций приведена на рисунке \ref{fig:full-volume-byX}.
% Схема расчета интегральной кривой от строевой по шпангоутам в табличном редакторе
\begin{figure}
    \centering
    \footnotesize
    \begin{tikzpicture}
        \matrix (mat) [table,
        row 1/.style={nodes={fill=gray!10,align=center,font=\segoe}},
        column 1/.style={nodes={fill=gray!10,align=center,font=\segoe}},
        column 1/.append style={nodes={text width=0.5cm}},
        column 2/.append style={nodes={text width=1.5cm}},
        column 3/.append style={nodes={text width=1.5cm}},
        column 4/.append style={nodes={text width=4.0cm}}
        ] {
        & A  & B & C \\
        1 & 0 & $\Omega_0$ & 0  \\
        2 & 1 & $\Omega_1$ & =C1+$\Delta L$*(B1+B2) \\
        3 & 2 & $\Omega_2$ & =C2+$\Delta L$*(B2+B3) \\
        4 & 3 & $\Omega_3$ & =C3+$\Delta L$*(B3+B4) \\
        5 & 4 & $\Omega_4$ & =C4+$\Delta L$*(B4+B5) \\
        6 & 5 & $\Omega_5$ & =C5+$\Delta L$*(B5+B6) \\
        7 & 6 & $\Omega_6$ & =C6+$\Delta L$*(B6+B7) \\
        8 & 7 & $\Omega_7$ & =C7+$\Delta L$*(B7+B8) \\
        9 & 8 & $\Omega_8$ & =C8+$\Delta L$*(B8+B9) \\
        10 & 9 & $\Omega_9$ & =C9+$\Delta L$*(B9+B10) \\
        11 & 10 & $\Omega_{10}$ & =C10+$\Delta L$*(B10+B11) \\
        };
    \end{tikzpicture}
    \caption{Схема расчета интегральной кривой от строевой по шпангоутам в табличном редакторе}
    \label{fig:full-volume-byX}
\end{figure}

Для построения кривой предельных длин отсеков задают несколько значений абсцисс ЦТ затопленных отсеков (обычно 5—7).

С кривой предельных объемов отсеков \ref{fig:max-volume-curve} для выбранных значений абсцисс ЦТ снимают значения объемов.
Полученные значения объемов в принятом масштабе в виде вертикальных отрезков наносят на интегральную кривую так, чтобы получить равенство площадей справа под и слева над интегральной кривой.
Для этого из нижней и верхней точек отрезка, соответствующего объему затопленного отсека, строятся горизонтальные прямы до пересечения с интегральной кривой.
Отрезок перемещается в вертикальном направлении так, чтобы площади совпали.

Горизонтальное расстояние между точками пересечения горизонтальных прямых, проведенных из нижней и верхней точки отрезка, соответствующего объему затопленного отсека, является предельной длиной отсека.

% К расчету кривой предельных длин отсеков
\begin{figure}[h!]
    \centerline{\incfig{figure_15 - 3}}
    \caption{К расчету кривой предельных длин отсеков}
    \caption*{\textit{
        Объем $v_i$ размещается таким образом, чтобы площади фигур $s_1$ и $s_2$ были равны.\\
        Расстояние между крайними точками площадей $l_i$ --- предельная длина отсека для положения $x^{\prime}_i$ по длине судна.
    }}
    \label{fig:max-length-curve}
\end{figure}

Полученные таким образом предельные длины отсеков для нескольких значений абсцисс ЦТ затопленных отсеков переносят на новую диаграмму.
По горизонтальной оси совпадающую с предыдущими диаграммами.

Длина предельного отсека откладывается вертикально вверх из точки, соответствующей середине этой длины.

Верхние концы отрезков соединяют плавной кривой --- кривой предельных длин отсеков.
В оконечностях эта кривая ограничена прямыми, проводимыми под углом к оси абсцисс $\beta=\ \mathrm{arctg}2$.

Объемы затопленных отсеков, использованные при построении кривой предельных длин отсеков, теоретические, то есть рассчитаны без учета объема, занимаемого конструкциями корпуса, устройствами, оборудованием и механизмами.
Фактически объем воды, поступивший в отсек, будет меньше теоретического объема.
Коэффициент проницаемости --- отношение фактического объема воды к теоретическому объему отсека.

Коэффициент проницаемости всегда меньше единицы, а значит, полученные по теоретическим объемам предельные длины отсеков могут быть увеличены на величину, обратную коэффициенту проницаемости $1 / \mu$.

% Пример кривой предельных длин отсеков
\begin{figure}[h!]
    \centerline{\incfig{figure_15 - 4}}
    \caption{Пример кривой предельных длин отсеков}
    \caption*{\textit{}}
    \label{fig:max-length-diagram}
\end{figure}

\subsection{Расчет посадки и начальной остойчивости поврежденного судна}

В настоящем задании предлагается выполнить расчет посадки и начальной остойчивости поврежденного судна методом постоянного водоизмещения, то есть исключением поврежденного отсека.
Используются допущение о прямобортности судна в пределах изменения осадки и метацентрические формулы остойчивости.

В курсовой работе в качестве исходных данных рекомендуется принимать следующее:

-- затапливаемый отсек сообщается с забортной водой;

-- затапливаемый отсек располагается между вторым и четвертым шпангоутами.

Для определения положения носовой и кормовой переборок затапливаемого отсека задаются положением его середины (например, на третьем шпангоуте).
С кривой предельных длин отсеков для выбранного положения снимают значение предельной длины отсека.
Длина затапливаемого отсека принимается равной предельной длине отсека.

В общем случае затопленный отсек можно представить в виде криволинейной пирамиды, основания которой представляют собой площади носовой и кормовой переборок этого отсека.
Ватерлиния отсека будет иметь форму трапеции.
Если площади носовой и кормовой переборок равны, то ватерлиния отсека будет иметь форму прямоугольника.

Для определения элементов затопленного отсека необходимо определить площади его носовой и кормовой переборок с использованием масштаба Бонжана.

С масштаба Бонжана снимаются значения площадей второго, третьего и четвертого шпангоутов по действующую ватерлинию.
По этим значениям строится участок строевой по шпангоутам, на котором наносят абсциссы носовой и кормовой переборок рассматриваемого отсека.

% Определение элементов затопленного отсека
\begin{figure}[h!]
    \centerline{\incfig{figure_16 - 1}}
    \caption{Определение элементов затопленного отсека}
    \caption*{\textit{
        $1$ --- кривая строевой по шпангоутам;
        $2$ --- кривая расчетной ватерлинии.
    }}
    \label{fig:comp-elements}
\end{figure}

Кривая строевой по шпангоутам от носовой до кормовой переборок заменяется прямым отрезком так, чтобы площади фигур были равны (см. рисунок \ref{fig:comp-elements}).
Значения площадей носовой и кормовой переборок принимаются равными основаниям получившейся трапеции (на рисунке --- $\Omega_{I}$ и $\Omega_{II}$).

Объем затопленного отсека может быть приближенно вычислен как объем пирамиды с основаниями $\Omega_{I}$ и $\Omega_{II}$ и высотой $l$

\begin{equation}
    v=\frac{\Omega_{I}+\Omega_{II}}{2}\cdot l.
\end{equation}

Объем влившейся воды определяется с учетом коэффициента проницаемости

\begin{equation}
    v_{\text{в}} = \mu \cdot v.
\end{equation}

Для более точного определения объема $v$ можно вычислить площадь строевой по шпангоутам между носовой и кормовой переборками отсека.

Координаты центра тяжести затопленного отсека можно определить как

\begin{equation}
    \begin{split}
        x_{\text{р}} = x_{II} + x_{\text{рc}}, \\
        y_{\text{р}} = 0, \\
        z_{\text{р}} = T - z_{WL},
    \end{split}
\end{equation}
где
$x_{\text{рc}}$ --- абсцисса центра тяжести отсека, отсчитываемая от кормовой его переборки;
$z_{WL}$ --- аппликата центра тяжести отсека, отсчитываемая от плоскости ватерлинии.

Величину $x_{\text{рc}}$ приближенно можно определить по формуле для центра тяжести пирамиды (на строевой по шпангоутам это прямоугольная трапеция)

\begin{equation}
    x_{\text{рc}} = \frac{1}{3} \cdot l \frac{\Omega_{II}+ 2 \Omega_{I}}{\Omega_{I}+\Omega_{II}}.
\end{equation}

Для вычисления $z_{WL}$ можно использовать формулу линейной аппроксимации

\begin{equation}
    z_{WL} = z_{WL}^{(3)} - \frac{z_{WL}^{(4)} - z_{WL}^{(2)}}{2 \Delta L} \cdot x_{\text{рC}},
\end{equation}
где
$z_{WL}^{(j)}$ --- отстояние центра тяжести j-го шпангоута от ватерлинии,
$j$ --- номер шпангоута (второй, третий, четвертый).

Для вычисления $z_{WL}^{(j)}$ используется формула

\begin{equation}
    z_{WL}^{(j)} = \frac{1}{\Omega^{(j)}} \int_0^T \Omega dz.
\end{equation}

Более точный расчет положения центра тяжести отсека можно выполнить с помощью строевой по шпангоутам, определяя статические моменты объема затопленного отсека относительно плоскости мидель-шпангоута ($YZ$) и основной плоскости ($XY$) по формулам

\begin{equation}
    \begin{split}
        x_{\text{р}} = \frac{1}{v} \int_{x_{II}}^{x_{I}} \Omega x dx, \\
        z_{\text{р}} = T - \frac{1}{v} \int_{x_{II}}^{x_{I}} \int_0^T \Omega dz dx.
    \end{split}
\end{equation}


\textit{Элементы потерянной площади ватерлинии}

Элементы потерянной площади ватерлинии определяются аналогично элементам потерянного объема по снятым с корпуса ординатам второго, третьего и четвертого шпангоутов.

Ординаты носовой и кормовой переборок на действующей ватерлинии определяют графическим методом с применением CAD–редактора.
Для этого на вспомогательном теоретическом чертеже (вид корпус) проводится действующая ватерлиния.
С проведенной ватерлинии на пересечении со вторым, третьим и четвертым шпангоутами снимаются значения ординат.

Ординаты на носовой и кормовой переборке можно определить графически с проведенной вспомогательной ватерлинии, или с применением линейной интерполяций по снятым значениям ординат второго, третьего и четвертого шпангоутов.

Площадь отсека определяется приближенно как площадь двух трапеций, построенных с использованием ординат кормовой переборки отсека, ординаты на третьем шпангоуте и ординаты носовой переборки отсека.


Потерянная площадь корректируется с учетом коэффициента заполнения $\mu_{S}$.

Абсцисса центра тяжести потерянной площади ватерлинии (площади затопленного отсека) определяется аналогично определению центра тяжести площади ватерлинии в Задании 1.

\begin{equation}
    x_{sc} = \frac{M_{sy}}{s_0}.
\end{equation}

Центральные моменты инерции потерянной площади ватерлинии относительно осей, параллельных осям $X$ и $Y$, находятся по формулам

\begin{equation}
    \begin{split}
        i_x = \mu_s \frac{l}{6} \frac{y_{II}^4 - y_{I}^4}{y_{II}-y_{I}}, \\
        i_y = \mu_s i_y^\prime - s \cdot(x_3 - x_s)^2,
    \end{split}
\end{equation}
где
$i_y^\prime$ -- момент инерции потерянной площади ватерлинии относительно третьего шпангоута;
$s \cdot(x_3 - x_s)^2$ --- переносный момент инерции к центру тяжести потерянной площади.

Для $i_y^\prime$ используется формула для трапеции

\begin{equation}
    i_y^\prime = l^3 \frac{3y_{II}+y_{I}}{24}.
\end{equation}

Имея исходные данные по судну и получив элементы затопленного отсека, вычисляют характеристики плавучести и остойчивости судна в поврежденном состоянии.

Изменение осадки определяется как

\begin{equation}
    \delta T = \frac{v_{\text{в}}}{S - s}.
\end{equation}

Координаты центра тяжести площади поврежденной ватерлинии определяются как

\begin{equation}
    \begin{split}
        x_f^\prime = \frac{S x_f - s x_s}{S - s}, \\
        y_f^\prime = - y_s \frac{s}{S - S}.
    \end{split}
\end{equation}

Моменты инерции $I_x^\prime$ и $I_{yf}^\prime$ площади поврежденной ватерлинии относительно центральных осей, параллельных осям OX и OY вычисляются как

\begin{equation}
    \begin{split}
        I_x^\prime = I_x - i_x - s y_s^2 (1 + \frac{s}{S - s}), \\
        I_{yf}^\prime = I_{yf} - i_y - s (x_s - x_F) (1 + \frac{s}{S - s}), \\
        I_{xf}^\prime = -s (x_s - x_f) y_s (1 + \frac{s}{S - s}).
    \end{split}
\end{equation}

Изменение координат центра величины вычисляется как

\begin{equation}
    \begin{split}
        \delta x_c = - \frac{v_{\text{в}}}{V} [x_{\text{р}} - x_f + (x_s - x_f)\frac{s}{S - s}], \\
        \delta z_c = - \frac{v_{\text{в}}}{V} [z_{\text{р}} - T - \frac{\delta T}{2}].
    \end{split}
\end{equation}

Аварийное возвышение центра величины вычисляется как

\begin{equation}
    z_c^\prime = z_c + \delta z_c.
\end{equation}

Аварийные метацентрические радиусы вычисляются как

\begin{equation}
    \begin{split}
        r^\prime = \frac{I_x^\prime}{V}, \\
        R^\prime = \frac{I_{yf}^\prime}{V}.
    \end{split}
\end{equation}

Аварийные метацентрические высоты

\begin{equation}
    \begin{split}
        h^\prime = z_c^\prime + r^\prime - z_g, \\
        H^\prime = z_c^\prime + R^\prime - z_g.
    \end{split}
\end{equation}

Угол дифферента по приближенной формуле А.Н.~Крылова

\begin{equation}
    \begin{split}
        \psi = \frac{v_{\text{в}} (x_{\text{р}} - x_f)}{V H^\prime} (1 + \frac{k_T}{k_x}), \\
        k_T = \frac{s}{S- - s}, \\
        k_x = \frac{x_{\text{р}} - x_F}{x_s - x_f}.
    \end{split}
\end{equation}

Аварийные осадки носом и кормой.

\begin{equation}
    \begin{split}
        T_{\text{Н}} = T + \delta T + \psi (L/2 - x_f^\prime), \\
        T_{\text{К}} = T + \delta T - \psi (L/2 + x_f^\prime).
    \end{split}
\end{equation}

В результате расчетов аварийной посадки и остойчивости делается вывод о том, соответствуют ли эти параметры требованиям правил проектирования и постройки.

\section{Продольный спуск}

Расчет продольного спуска состоит из двух частей:

-- расчета и построения английской диаграммы спуска;

-- определения наибольшего погружения заднего конца полоза.


Диаграмма строится для второго периода продольного спуска, когда центр тяжести судна движется параллельно линии спусковых дорожек.
Диаграмма позволяет определить:

-- наличие и момент опрокидывания;

-- начало всплытия (начало третьего периода);

-- наибольшее значение баксова давления;

-- величину наибольшего погружения переднего конца полоза.

Статический расчет второго периода продольного спуска выполняется с использованием масштаба Бонжана.

Основные исходные данные для расчета продольного спуска, кроме кривых элементов теоретического чертежа и масштаба Бонжана, являются следующие величины:

-- спусковой вес судна $D_{\text{С}}$ (рекомендуется принимать его равным четверти водоизмещения судна по КВЛ, $D_{\text{С}} = 0{,}25 D_{\text{КВЛ}}$);

-- абсцисса и аппликата центра тяжести судна на момент спуска, $x_g$, $z_g$;

-- длина полоза L (рекомендуется принимать равной $0{,}8L_{\text{КВЛ}}$).
Середина полоза совпадает с положением мидель-шпангоута.
Тогда длина переднего конца полоза $L_1 = 0{,}5 L + x_g$, заднего конца полоза $L_2 = 0{,}5 L - x_g$;

-- расчетная глубина воды на пороге $T_0$ (рекомендуется принимать равной 0,8 от осадки спущенного на воду судна);

-- возвышение киля над зеркалом спусковых дорожек принимается равным для переднего конца полоза $C_1=0{,}70\mathrm{м}$, для заднего конца полоза $C_2=0{,}50\mathrm{м}$;

-- суммарная ширина полозьев принимается равной 2 м;

-- уклон спусковых дорожек принимается равным $\beta=\ 0{,}06$;

-- плотность воды в районе спуска принимается равной $\rho=1{,}000\mathrm{т/м}^3$.

% Схема спускового устройства
\begin{figure}[h!]
    \centerline{\incfig{figure_17 - 1}}
    \caption{Схема спускового устройства}
    \caption*{\textit{}}
    \label{fig:launching-device}
\end{figure}


Далее приводится последовательность выполнения расчетов продольного спуска.

Вычисляют уклон линии киля

\begin{equation}
    \alpha = \beta + \frac{C_2 - C_1}{L_1 + L_2}.
\end{equation}

Длина подводной части дорожек вычисляется по формуле

\begin{equation}
    \lambda = \frac{T_0}{\beta}.
\end{equation}

С кривым элементов теоретического чертежа для спускового веса судна определяют его среднюю осадку, координаты центра величины, абсциссу центра тяжести площади ватерлинии, продольный метацентрический радиус.

Вычисляют продольную метацентрическую высоту

\begin{equation}
    H = R + z_c - z_g.
\end{equation}

С использованием полученного значения продольной начальной метацентрической высоты определяют дифферент судна

\begin{equation}
    \psi = \frac{x_g - x_c}{H}.
\end{equation}

Тогда углубление передним и задним концами полоза будут определяться по формулам

\begin{equation}
    \begin{split}
        T_1 = T + C_1 - \psi (L_1 + x_f^\prime), \\
        T_2 = T + C_2 + \psi (L_2 - x_f^\prime),
    \end{split}
\end{equation}
где
$x_f^\prime = x_f - x_g$ --- отстояние центра тяжести площади ватерлинии от центра тяжести судна.

\textit{Контрольный шаг:} если $T_2 > T_0$, то сход судна со стапеля (в четвертом периоде) будет сопровождаться соскакиванием.

Баксово давление может быть определено по приближенной формуле

\begin{equation}
    N_{\text{Б}} = \frac{D_{\text{С}} H} {L_2 - x_f^\prime} (\alpha + \psi).
\end{equation}

С учетом полученного значения баксова давления определяется водоизмещение всплытия

\begin{equation}
    V_2 = \frac{1}{\rho} (D_{\text{С}} - N_{\text{Б}}).
\end{equation}

Для полученного значения водоизмещения $V_2$ определяется значение средней осадки $T_{\text{СР}}^\prime$ и положение центра тяжести действующей ватерлинии $x_{f2}$.

Приближенная формула высоты ватерлинии всплытия (с учетом запаса 20\% )

\begin{equation}
    h = 1{,}2 (T_{\text{СР}}^\prime + x_{f2}^\prime \alpha),
\end{equation}
где
$x_{f2}^\prime = x_{f2} - x_g$.

Определяют осадки на носовом и кормовом перпендикулярах, соответствующие ватерлинии всплытия

\begin{equation}
    \begin{split}
        T_{\text{Н}} = h - (0,5 L - x_g) \cdot \alpha, \\
        T_{\text{К}}  = h + (0,5 L + x_g) \cdot \alpha.
    \end{split}
\end{equation}

Для выполнения дальнейших этапов расчета продольного спуска выбирается количество расчетных ватерлиний $n$ в диапазоне от пяти до восьми.

Шаг между расчетными ватерлиниями определяется как $\Delta h = h / n$ и округляется с точностью до 5 см.

Расчетные ватерлинии вплоть до принятой ватерлинии всплытия наносятся на масштаб Бонжана.

% Масштаб Бонжана с нанесенными расчетными ватерлиниями
\begin{figure}[h!]
    \centerline{\incfig{figure_17 - 2}}
    \caption{Масштаб Бонжана с нанесенными расчетными ватерлиниями}
    \caption*{\textit{На масштабе Бонжана нанесены расчетные ватерлинии, а также ватерлиния всплытия.}}
    \label{fig:bonjan-scale}
\end{figure}

Для каждой расчетной ватерлинии выполняется расчет водоизмещения и статического момента $M$ относительно плоскости $YZ$ и $M^\prime$ относительно плоскости $Y^\prime Z^\prime$

\begin{equation}
    M^\prime = M - V \cdot x_g.
\end{equation}

Английская диаграмма в зависимости отражает изменение следующих величин от пути $S$, пройденного судном во втором периоде:

-- силы плавучести $\rho W$;

-- силы веса $D_{\text{С}}$ --- прямая горизонтальная линия;

-- момента сил плавучести относительно заднего конца полоза $M_W$;

-- момента силы веса относительно заднего конца полоза ($M_D = D_{\text{С}} L_2$) --- прямая горизонтальная линия;

--- момента силы плавучести относительно порога $M_W^\prime$;

--- момента силы веса относительно порога $M_D^\prime = D_{\text{С}} a$ ($a$ --- расстояние от центра тяжести до порога) --- наклонная прямая.

Расчет силы плавучести и моментов силы плавучести выполняют в табличной форме.
При этом используются следующие выражения:

Объемное водоизмещение погруженной части судна определяется как

\begin{equation}
    W = V - v^\prime.
\end{equation}

Потерянный объем плавучести определяется по формуле

\begin{equation}
    v^\prime = 1/2 \beta l_{\text{п}}^2,
\end{equation}
где 
$l_{\text{п}}$ --- длина спускового полоза, вошедшего в воду и прилегающего к спусковым дорожкам, равная меньшей из величин $\lambda$ или $S$ (см. рисунок).


По рассчитанным величинам выполняют построение английской диаграммы спуска.

По оси абсцисс в масштабе откладывается путь, пройденный судном во втором периоде, по оси ординат (вертикальной оси) строятся две шкалы --- для сил и для моментов.
Диапазон шкал следует выбирать таким, чтобы кривые изменения силы плавучести и ее моментов с одной стороны не сливались на диаграмме, а с другой стороны полностью отображались на ней, не выходя за границы.

По полученным в расчетах величинам силы плавучести и ее моментов от пройденного пути на диаграмму наносятся соответствующие кривые.

На диаграмму наносятся две горизонтальные линии, соответствующие силе веса и ее моменту относительно заднего конца полоза.
Также наносится наклонная прямая момента сил веса относительно порога $M_D^\prime = D_{\text{С}} (S - \lambda - L_1)$.

Затем выполняют анализ построенной диаграммы.

% Пример английской диаграммы спуска
\begin{figure}[h!]
    \centerline{\incfig{figure_17 - 3}}
    \caption{Пример английской диаграммы спуска}
    \caption*{\textit{1 --- центр тяжести судна на пороге;
    2 --- критическое положение судна;
    3 --- момент всплытия.}}
    \label{fig:english-diagram}
\end{figure}

Если момент сил плавучести больше момента сил тяжести относительно порога, это значит, что опрокидывание отсутствует.

Определяется критическое положение судна, при котором разница между моментами сил плавучести и сил тяжести относительно порога минимальна.

Определяется расстояние, при котором происходит всплытие.
Условие --- равенство моментов силы тяжести и плавучести относительно заднего конца полоза $M_W = M_D$.

Определяется наибольшее баксово давление в момент всплытия $N_{\text{Б}} = D_{\text{С}} - \rho W$.

Наибольшее погружение переднего конца полоза определяется по формуле

\begin{equation}
    T_1 = S \cdot \beta.
\end{equation}

Величина наибольшего погружения заднего конца полоза после соскакивания определяется по формуле

\begin{equation}
    T_{2 \text{MAX}} = T_2 + \Delta T.
\end{equation}

Переуглубление заднего конца полоза из-за прыжка в первом приближении (пренебрегая сопротивлением воды) можно определить по формуле

\begin{equation}
    \Delta T = T_{2} = T_0.
\end{equation}
